\documentclass[11pt,a4paper]{scrartcl} 


\usepackage{../manostilius_lt}

% Projekto informacija (būtina preambulėje dėl jūsų stiliaus specifikos)
\title{Laipsnis su sveikuoju rodikliu}
\author{Andrius Berniukevičius} 

\begin{document}

\maketitle


\section{Laipsnio su sveikuoju rodikliu apibrėžimas}
Norint išplėsti laipsnio su natūraliuoju rodikliu apibrėžtį iki sveikųjų skaičių, būtina nustatyti laipsnio reikšmes, kai rodiklis yra lygus nuliui arba neigiamam sveikajam skaičiui. 
\begin{proposition}[Laipsnio, kurio rodiklis lygus nuliui, reikšmė visada lygi $1$ ]
Jei $a \neq 0$, tai $a^0 = 1$.
\end{proposition}
\begin{proof}
Laipsnių su vienodais pagrindais dalybos savybė teigia, kad bet kokiems natūraliesiems skaičiams $m$ ir $n$ (kai $m > n$) galioja savybė:
$$\frac{a^m}{a^n} = a^{m-n}$$
Sukurkime situaciją, kai laipsnio rodikliai yra lygūs, t. y. $m = n$. Užrašykime dviejų vienodų laipsnių dalybą $\frac{a^n}{a^n}$ ir išnagrinėkime ją dviem būdais:
\begin{enumerate}
    \item \textbf{Remiantis aritmetikos taisykle:} Bet kokį nelygų nuliui skaičių ar reiškinį padaliję iš savęs, gauname vienetą:
    $$\frac{a^n}{a^n} = 1$$
    \item \textbf{Remiantis laipsnių dalybos savybe:} Atimame rodiklius:
    $$\frac{a^n}{a^n} = a^{n-n} = a^0$$
\end{enumerate}
Kadangi abiem būdais buvo pertvarkomas tas pats reiškinys $\frac{a^n}{a^n}$, gautieji rezultatai privalo būti lygūs:
$$a^0 = 1$$
\end{proof}

\begin{proposition} [Laipsnis su neigiamu sveikuoju rodikliu]
Jei $a \neq 0$ ir $k$ yra natūralusis skaičius, tai $a^{-k} = \frac{1}{a^k}$.
\end{proposition}
\begin{proof}
Paimkime bet kokį natūralųjį skaičių $m$. Kadangi $k$ taip pat yra natūralusis skaičius, sudarykime laipsnių su vienodais pagrindais dalybos reiškinį, kuriame vardiklio rodiklis yra didesnis už skaitiklio rodiklį tiksliai skaičiumi $k$:
$$\frac{a^m}{a^{m+k}}$$
Išnagrinėkime šį reiškinį dviem skirtingais būdais, remdamiesi tik natūraliųjų laipsnių savybėmis:
\begin{enumerate}
    \item \textbf{Prastinimas pagal apibrėžimą:} \\
    Pagal laipsnių daugybos savybę, vardiklį galime išskleisti kaip laipsnių su natūraliaisiais rodikliais sandaugą: $a^{m+k} = a^m \cdot a^k$. Įrašykime tai į trupmeną:
    $$\frac{a^m}{a^{m+k}} = \frac{a^m}{a^m \cdot a^k}$$
    Kadangi skaitiklyje ir vardiklyje turime vienodą natūralųjį laipsnį $a^m$, jį galime suprastinti:
    $$\frac{\cancelto{1}{a^m}}{\cancelto{1}{a^m} \cdot a^k} = \frac{1}{a^k}$$
    \item \textbf{Laipsnių savybių taikymas:} \\
    Jeigu pritaikytume  laipsnių dalybos savybę $\frac{a^m}{a^n} = a^{m-n}$ tiesiogiai, gautume:
    $$\frac{a^m}{a^{m+k}} = a^{m - (m + k)} = a^{m - m - k} = a^{-k}$$
\end{enumerate}
Kadangi abu sprendimo būdai gauti pertvarkant tą patį pradinį reiškinį, jų galutiniai rezultatai privalo būti lygūs:
$$a^{-k} = \frac{1}{a^k}$$
\end{proof}
\begin{corollary}
Keliant trupmeną neigiamu laipsniu, ji yra apverčiama, o laipsnio rodiklis tampa teigiamas:
$$\left(\frac{a}{b}\right)^{-k} = \left(\frac{b}{a}\right)^k$$
\end{corollary}
\begin{proof}
Šį teiginį įrodysime kaip išvadą iš jau įrodytos lygybės 
$$x^{-k} = \frac{1}{x^k}$$
Tarkime, kad mūsų laipsnio pagrindas yra trupmena $x = \frac{a}{b}$. Pritaikykime jai neigiamo laipsnio apibrėžimą:
$$\left(\frac{a}{b}\right)^{-k} = \frac{1}{\left(\frac{a}{b}\right)^k}$$
Vardikliui pritaikome trupmenos kėlimo laipsniu savybę:
$$\frac{1}{\left(\frac{a}{b}\right)^k} = \frac{1}{\frac{a^k}{b^k}}$$
Gautą užrašą suprantame kaip vieneto dalybą iš trupmenos $\frac{a^k}{b^k}$. Žinome, kad dalijant iš trupmenos, veiksmas pakeičiamas daugyba iš jai atvirkštinės trupmenos:
$$1 : \frac{a^k}{b^k} = 1 \cdot \frac{b^k}{a^k} = \frac{b^k}{a^k}$$
Gautai trupmenai $\frac{b^k}{a^k}$ pritaikome laipsnių savybę atvirkštine tvarka (iškeliame bendrą laipsnio rodiklį $k$ už skliaustų):
$$\frac{b^k}{a^k} = \left(\frac{b}{a}\right)^k$$
Sujungę visą mąstymo grandinę nuo pradžios iki galo, gauname:
$$\left(\frac{a}{b}\right)^{-k} = \left(\frac{b}{a}\right)^k$$
\end{proof}
Dabar galime suformuluoti laipsnio su sveikuoju rodikliu apibrėžimą:
\begin{definition}
   \vocab{Laipsnis su sveikuoju rodikliu} yra laipsnis $a^n$, kur  $a \neq 0$ ir $n$ yra sveikasis skaičius.
\end{definition}
Galimi trys atvejai:
\begin{enumerate}
    \item \textbf{Kai rodiklis teigiamas ($n > 0$):} 
    $$a^n = \underbrace{a \cdot a \cdot \dots \cdot a}_{n \text{ kartų}}$$
    
    \item \textbf{Kai rodiklis lygus nuliui ($n = 0$):} 
    $$a^0 = 1$$
    
    \item \textbf{Kai rodiklis neigiamas ($n < 0$, t. y. $n = -k$, kur $k$ -- natūralusis skaičius):}
    $$a^{-k} = \frac{1}{a^k}$$
\end{enumerate}
\begin{remark}
Jei pagrindas yra paprastoji trupmena, neigiamas sveikasis rodiklis ją tiesiog apverčia:

$$\left(\frac{a}{b}\right)^{-k} = \left(\frac{b}{a}\right)^k $$
\end{remark}

Jei pagrindas yra paprastoji trupmena, neigiamas sveikasis rodiklis ją tiesiog apverčia:

$$\left(\frac{a}{b}\right)^{-k} = \left(\frac{b}{a}\right)^k = \frac{b^k}{a^k}$$

\section{Uždavinių sprendimo pavyzdžiai}

\begin{example}
Išspręskite lygtį:
$$\left(\frac{2}{5}\right)^x = 6\frac{1}{4}$$\\
\textbf{Sprendimas:} Mišrųjį skaičių užrašome netaisyklingąja trupmena ir abi lygties puses išreiškiame tuo pačiu pagrindu:
$$\left(\frac{2}{5}\right)^x = \frac{25}{4}$$
$$\left(\frac{2}{5}\right)^x = \left(\frac{5}{2}\right)^2 = \left(\frac{2}{5}\right)^{-2}$$
Kadangi laipsnių pagrindai vienodi, jų rodikliai taip pat lygūs:
$$x = -2$$\\
\textbf{Atsakymas:} $x = -2$.
\end{example}

\begin{example}
Supaprastinkite reiškinį:
$$\frac{28a^6b^4c^{-13}}{24a^{10}b^{-3}c}$$\\
\textbf{Sprendimas:} Taikome laipsnių dalybos savybę ir neigiamus rodiklius pertvarkome į teigiamus:
\[
\begin{aligned}
\frac{28a^6b^4c^{-13}}{24a^{10}b^{-3}c}
&= \frac{28}{24} \cdot \frac{a^6}{a^{10}} \cdot \frac{b^4}{b^{-3}} \cdot \frac{c^{-13}}{c^1} \\
&= \frac{7}{6} \cdot a^{-4} \cdot b^7 \cdot c^{-14} \\
&= \frac{7}{6} \cdot \frac{1}{a^4} \cdot b^7 \cdot \frac{1}{c^{14}} \\
&= \frac{7b^7}{6a^4c^{14}}.
\end{aligned}
\]\\
\textbf{Atsakymas:} $\dfrac{7b^7}{6a^4c^{14}}$.
\end{example}

\begin{example}
Supaprastinkite reiškinį:
$$\left(\frac{3}{10}x^{-4}yz^2\right)^{-4}$$

\textbf{Sprendimas:} Pirmiausia neigiamą rodiklį pertvarkome į teigiamą:
$$\left(\frac{3}{10}x^{-4}yz^2\right)^{-4} = \left(\frac{3yz^2}{10x^4}\right)^{-4}$$
Kadangi trupmena keliama neigiamu laipsniu, ją apverčiame:
$$\left(\frac{3yz^2}{10x^4}\right)^{-4} = \left(\frac{10x^4}{3yz^2}\right)^4$$
Keliame skaitiklį ir vardiklį ketvirtuoju laipsniu:
$$\left(\frac{10x^4}{3yz^2}\right)^4 = \frac{10^4 \cdot (x^4)^4}{3^4 \cdot y^4 \cdot (z^2)^4} = \frac{10000x^{16}}{81y^4z^8}$$

\textbf{Atsakymas:} $\dfrac{10000x^{16}}{81y^4z^8}$.
\end{example}

\begin{example}
Supaprastinkite reiškinį:
$$\left(\frac{6x^{-2}y^3}{5ab^4}\right)^{-2} : \frac{5a^{-3}b}{72xy^4}$$

\textbf{Sprendimas:} Pirmiausia neigiamus rodiklius pertvarkome į teigiamus:
$$\left(\frac{6x^{-2}y^3}{5ab^4}\right)^{-2} : \frac{5a^{-3}b}{72xy^4} = \left(\frac{6y^3}{5ab^4x^2}\right)^{-2} : \frac{5b}{72a^3xy^4}$$
Pirmąją trupmeną apverčiame ir keliame antruoju laipsniu:
$$\left(\frac{6y^3}{5ab^4x^2}\right)^{-2} = \left(\frac{5ab^4x^2}{6y^3}\right)^2 = \frac{25a^2b^8x^4}{36y^6}$$
Dalybą iš antrosios trupmenos pakeičiame daugyba iš jai atvirkštinės trupmenos:
$$\frac{25a^2b^8x^4}{36y^6} \cdot \frac{72a^3xy^4}{5b} = 10a^5b^7x^5y^{-2} = \frac{10a^5b^7x^5}{y^2}$$

\textbf{Atsakymas:} $\dfrac{10a^5b^7x^5}{y^2}$.
\end{example}

\section{Uždaviniai}
\begin{problem}
Atlikite veiksmus su laipsniais. Galutiniuose atsakymuose laipsnius palikite tik su teigiamais rodikliais:

\begin{tasks}[label={(\arabic*)}, label-offset=0.9em](2)
    \task $6^x=\frac{1}{216}$
    \task $\left(\frac{3}{4}\right)^x=2\frac{10}{27}$
    \task $0,2^x=125$
    \task $\left(\frac{3}{4}a^3b^{-2}\right)^{-2}$
    \task $\left(\frac{2x^3}{3y^4}\right)^{-4}$
    \task $\left(\frac{2m^3}{5n^7}\right)^{-2}\cdot\frac{4m^3}{125n^{10}}$
    \task $2,5^x=\frac{8}{125}$
    \task $5^x=0,0016$
    \task $\left(0,1x^5y^{-4}z\right)^{-3}$
    \task $\frac{2x^3y^{-6}z}{12x^4y^6z^{-3}}$
    \task $\frac{36a^5bc^{-10}d^3}{24ab^{-1}c^{-3}d^6}$
    \task $\frac{45mn^{10}p^{-3}}{25m^{-3}n^{1-1}p^4}$
    \task $\left(\frac{3ab^4}{5xy}\right)^{-1}\cdot\frac{9x^{-1}y^2}{10ab^4}$
    \task $4^x=\frac{1}{64}$
    \task $\left(\frac{2}{3}\right)^x=\frac{27}{8}$
    \task $0,25^x=64$
    \task $\left(\frac{5}{2}c^{-2}d^3\right)^{-2}$
    \task $\left(\frac{3p^2}{4q^3}\right)^{-3}$
    \task $\left(\frac{7r^{-3}s^2}{2t}\right)^{-2}\cdot\frac{49r^2}{4s^4}$
    \task $0,4^x=\frac{125}{8}$
    \task $\left(\frac{1}{3}\right)^x=81$
    \task $\left(0,2u^{-3}v^2w^{-1}\right)^{-2}$
    \task $\left(\frac{4a^{-2}b^3}{9c}\right)^{-1} : \frac{27ac^2}{8b^5}$
\end{tasks}
\end{problem}

\newpage
\section{Atsakymai}
\begin{tasks}[label={(\arabic*)}, label-offset=0.9em](2)
    \task $x=-3$
    \task $x=-3$
    \task $x=-3$
    \task $\dfrac{16b^4}{9a^6}$
    \task $\dfrac{81y^{16}}{16x^{12}}$
    \task $\dfrac{n^4}{5m^3}$
    \task $x=-3$
    \task $x=-4$
    \task $\dfrac{1000y^{12}}{x^{15}z^3}$
    \task $\dfrac{z^4}{6xy^{12}}$
    \task $\dfrac{3a^4b^2}{2c^7d^3}$
    \task $\dfrac{9m^4n^{10}}{5p^7}$
    \task $\dfrac{3y^3}{2a^2b^8}$
    \task $x=-3$
    \task $x=-3$
    \task $x=-3$
    \task $\dfrac{4c^4}{25d^6}$
    \task $\dfrac{64q^9}{27p^6}$
    \task $\dfrac{r^8t^2}{s^8}$
    \task $x=-3$
    \task $x=-4$
    \task $\dfrac{25u^6w^2}{v^4}$
    \task $\dfrac{2ab^2}{3c}$
\end{tasks}





\end{document}